\documentclass[10pt,a4paper]{amsart}
\usepackage[margin=19mm]{geometry}
\usepackage{amsmath,amssymb,mathtools}
\usepackage[scr=rsfs,cal=euler]{mathalfa}
\usepackage{xfrac}
\usepackage[unicode,hidelinks,bookmarks=false]{hyperref}
\newcommand{\ie}{i.e.\ }
\newcommand{\cf}{cf.\ }
\newcommand{\resp}{respectively\ }
\newcommand{\Subsection}{\subsection}
\providecommand{\nobreakdash}{\nobreak\hbox{-}}
\setlength{\emergencystretch}{2em}
\multlinegap0pt
\usepackage{xcolor}
\usepackage{tcolorbox}
\usepackage{amssymb}
\usepackage{mathtools}
\usepackage[shortlabels]{enumitem}
\usepackage{tikz}
\newtheorem{proposition}{Proposition}
\newcommand{\RB}{\mathbb{R}}
\newcommand{\bM}{\overline{M}}
\newcommand{\bW}{\smash{\overline{W}}\vphantom{W}}
\newcommand{\bw}{\overline{w}}
\title{Dynamics of Gradient Descent with Large Step Size Near a Manifold of Flat Minima}
\author{Lachlan Ewen MacDonald, Ren\'e Vidal}
\date{Source: arXiv:2607.08380v2 (CC BY 4.0)}
\begin{document}
\maketitle

\noindent\emph{Source and licence.} Dynamics of Gradient Descent with Large Step Size Near a Manifold of Flat Minima. Version arXiv:2607.08380v2 (CC BY 4.0), distributed by arXiv under the Creative Commons Attribution 4.0 International licence, http://creativecommons.org/licenses/by/4.0/, as recorded in the arXiv licence metadata for this exact version and shown on the abstract page https://arxiv.org/abs/2607.08380v2. Full licence notice: \texttt{licenses/arxiv-2607.08380-CC-BY-4.0.txt}.\par\smallskip

\noindent\textit{Editorial scope.} Counted unit: the article's proposition proposition (source0.tex lines 2166--2176). The article's complete original proof of it is delivered with it (source0.tex lines 2178--2294, one \texttt{proof} environment of 117 lines). No mathematics is written, repaired or re-derived: the statement and the proof are the article's own lines, in the article's own notation, and no retained line is substituted or re-flowed.  Retained uncounted as the source-local context the proof needs: lines 1511--1513; lines 1583--1585; lines 1661--1693; lines 1738--1762; lines 1770--1772; lines 2134--2136; lines 2154--2156; lines 2157--2159; lines 2161--2163.  Nothing was omitted from any retained span, so the package carries no omission record.  No retained line contains a citation, so the wrapper carries no bibliography and no bibliography entry.  No retained line contains a figure environment, an image-inclusion command or drawing code, so no image is delivered and the package declares no asset dependency.  Editorial formatting only, in addition: an AMS \texttt{amsart} preamble that restates the article's own declarations and macros used by the retained lines, namely the environments \texttt{proposition} on separate counters, together with the standard packages those lines need.  The retained environments print in this document as Proposition 1 in order of appearance, whereas the article prints the same items with the numbering of its own full document.  The complete native source of the article as distributed in the arXiv e-print for this exact version is bundled unchanged as \texttt{sources/204775/source0.tex}.
\begin{align}\label{eq:matrixfactorisationf}
f(w):=W_L\cdots W_1,\qquad w:=(W_1,\dots,W_L)\in \prod_{l=1}^{L}\RB^{d_l\times d_{l-1}}.
\end{align}

    \begin{align}\label{eq:wnormalform}
    \Bigg(\begin{pmatrix} \sigma_1^{1/L} & 0 \\ 0 & \bW_l\end{pmatrix}\Bigg)_{l=1}^L
    \end{align}

\begin{proposition}\label{prop:tangents}
    At a point $w\in F$ of the form
    \begin{align}
    \Bigg(\begin{pmatrix} \sigma_1^{1/L} & 0 \\ 0 & \bW_l\end{pmatrix}\Bigg)_{l=1}^L
    \end{align}
    as in \eqref{eq:wnormalform}, with $\bar{w}=(\bW_1,\dots,\bW_L)\in\overline{M}$, given $a:=(a_l)_{l=1}^{L-1}\in \RB^{L-1}$, $b:=(b_l)_{l=1}^{L-1},c:=(c_l)_{l=1}^{L-1}\in\prod_{l=1}^{L-1}\RB^{d_l-1}$ and $D:=(D_l)_{l=1}^{L}\in T_{\bw}\bM$, denote
    \begin{align}
        v_1(w;a):=\Bigg(\begin{pmatrix} \sigma_1^{1/L}(a_l-a_{l-1}) & 0 \\ 0 & 0\end{pmatrix}\Bigg)_{l=1}^{L},
    \end{align}
    \begin{align}
        v_2(w;b,c):=\Bigg(\begin{pmatrix} 0 & \big(\bW_l^Tb_l - \sigma_1^{1/L}b_{l-1}\big)^T \\ \sigma_1^{1/L}c_l-\bW_lc_{l-1} & 0\end{pmatrix}\Bigg)_{l=1}^{L}
    \end{align}
    and
    \begin{align}
        v_3(w;D):=\bigg(\begin{pmatrix} 0 & 0 \\ 0 & D_l\end{pmatrix}\bigg)_{l=1}^{L},
    \end{align}
    where $a_0,b_0,c_0$ and $a_L,b_L,c_L$ are taken to be zero. Then all such $v_1(w;a),v_2(w;b,c),v_3(w;D)$ span the tangent space $T_wM$. Moreover, the tangent space to $F$ at $w$ is
    \begin{equation}
        T_wF = \mathrm{span}\bigg\{v_2(w,\xi,-\xi),v_3(w,D):\xi\in\prod_{l=1}^{L-1}\RB^{d_l-1},\,D\in T_{\bw}\bM\bigg\}
    \end{equation}
    and its orthogonal complement $\nu_wF$ in $T_wM$ is
    \begin{equation}
        \nu_wF = \mathrm{span}\bigg\{v_1(w,a),v_2(w,b,c):a\in\RB^{L-1},\,b,c\in\prod_{l=1}^{L-1}\RB^{d_l-1}\text{ s.t. }K_1(w)[b] = K_2(w)[c]\bigg\},
    \end{equation}
    where
    \begin{equation}
        K_1(w)[b]:=\bigg(\big(\sigma_1^{2/L}I_{d_l-1}+\bW_l\bW_l^T\big)b_l - \sigma_1^{1/L}\bW_lb_{l-1} - \sigma_1^{1/L}\bW_{l+1}^Tb_{l+1}\bigg)_{l=1}^{L-1}
    \end{equation}
    and
    \begin{equation}
        K_2(w)[c]:=\bigg(\big(\sigma_1^{2/L}I_{d_l-1}+\bW_{l+1}^T\bW_{l+1}\big)c_l-\sigma_1^{1/L}\bW_lc_{l-1}-\sigma_1^{1/L}\bW_{l+1}^Tc_{l+1}\bigg)_{l=1}^{L-1}.
    \end{equation}
\end{proposition}

\begin{proposition}\label{prop:morsebott}
    Let $w\in F$ be of the form \eqref{eq:wnormalform}. With $\lambda_1:=\lambda_1(Df^TDf)$, the subspaces
    \begin{align}
    V_1(w):=\mathrm{span}\bigg\{v_1(w;a):a\in\RB^{L-1}\bigg\},
    \end{align}
    \begin{align}
    V_2(w):=\mathrm{span}\bigg\{v_2(w;b,c):b,c\in\prod_{l=1}^{L-1}\RB^{d_l-1}\text{ s.t. }K_1(w)[b] = K_2(w)[c]\bigg\}
    \end{align}
    of $\nu_wF$ are invariant under $\nabla^2_M\lambda_1(w)$, with
    \begin{align}
    \mathrm{spec}(\nabla^2_M\lambda_1(w)|_{V_1(w)}) = 4\sigma_1^{2-4/L},
    \end{align}
    \begin{align}
    \mathrm{spec}(\nabla^2_M\lambda_1(w)|_{V_2(w)}) = \text{solutions $\mu(w)$ of generalised eigenproblem }H(w)v = \mu(w) G(w)v,
    \end{align}
    where, suppressing evaluation at $w$,
    \begin{align}
    H:=2\sigma_1^{2-2/L}\big(I_{\prod_{l=1}^{L-1}\RB^{d_l-1}} + \Gamma_{01}^T(\lambda_1I_{d_0-1}-\Delta_{01})^{-1}\Gamma_{01} + \Gamma_{10}^T(\lambda_1I_{d_L-1}-\Delta_{10})^{-1}\Gamma_{10}\big)
    \end{align}
    and
    \begin{align}
    G:=(K_1^{-1}+K_2^{-1})^{-1}
    \end{align}
    are positive-definite. Consequently, $\lambda_1$ is Morse-Bott along $F$.
\end{proposition}

    \begin{align}\label{eq:Theta}
    \Theta: (W_1,\dots,W_L) \mapsto \big(J_1W_1J_0,\dots,J_LW_LJ_{L-1}\big)
    \end{align}

\begin{align}\label{eq:alpha}
\alpha_{\eta}(x):=-\bigg(D^3\ell[\nu_1^{\otimes 2},A_{\eta}\nabla^3\ell[\nu_1^{\otimes 2}]](x) + \frac{1}{3}D^4\ell[\nu_1^{\otimes 4}](x)\bigg)
\end{align}

\begin{align}\label{eq:D3ell}
    D^3\ell = 3\langle D^2f\odot Df\rangle + \langle D^3f,f-\tau\rangle\Rightarrow D^3\ell|_M = 3\langle D^2f\odot Df\rangle,
\end{align}

\begin{align}\label{eq:D4ell}
    D^4\ell = 4\langle D^3f\odot Df\rangle + 3\langle D^2f\odot D^2f\rangle + \langle D^4f,f-\tau\rangle\Rightarrow D^4\ell|_M = 4\langle D^3f\odot Df\rangle + 3\langle D^2f\odot D^2f\rangle.
\end{align}

\begin{align}\label{eq:Dkf}
    D^kf(w)[\xi^1,\dots,\xi^k] = \sum_{1\leq l_1<\dots<l_k\leq L}\sum_{\pi\in\Pi(k)}W_{L:l_k+1}\xi^{\pi(1)}_{l_k}W_{l_k-1:l_{k-1}+1}\cdots W_{l_2-1:l_1+1}\xi^{\pi(k)}_{l_1}W_{l_{1}-1:1}.
\end{align}

\begin{proposition}
    For the matrix factorisation model $f$ of \eqref{eq:matrixfactorisationf}, the function $\alpha_{\eta}$ of \eqref{eq:alpha} is constant on $F$, with
    \begin{align}
    \alpha_{\eta}|_{F} = -\bigg(\frac{9(L-1)^2}{(1-\eta\lambda_1|_{F})} + \frac{(L-1)(7L-11)}{3}\bigg)\sigma_1^{2-4/L},
    \end{align}
    for all $\eta$ in a neighbourhood of $2/\lambda_1|_{F}$, and in particular satisfies
    \begin{align}
    \alpha_{2/\lambda_1}|_{F} = (L-1)\sigma_1^{2-4/L}\bigg(\frac{20L-16}{3}\bigg) >0
    \end{align}
    for all $L\geq 2$. Moreover, $D\alpha_{\eta}|_{F} = 0$.
\end{proposition}

\begin{proof}
    Since $\alpha_{\eta}$ is defined solely in terms of derivatives of $f$ and in terms of the metric, it suffices to do the computations at a point $w$ of the form \eqref{eq:wnormalform}, i.e.
    \begin{align}
    w  = \Bigg(\begin{pmatrix}\sigma_1^{1/L} & 0 \\ 0 & \bW_{l}\end{pmatrix}\Bigg)_{l=1}^L
    \end{align}
    for some $(\bW_1,\dots,\bW_L)\in \bM$. It is then clear that one has
    \begin{align}
    \nu_1(w) = \frac{1}{\sqrt{L}}\big(e_1(d_l)e_1(d_{l-1})^T\big)_{l=1}^L.
    \end{align}
    We begin by proving that $\alpha_{\eta}$ is constant on $F$ for all $\eta$ sufficiently close to $2/\lambda_1|_{F}$. Suppressing evaluation at $w$ for notational convenience, using \eqref{eq:Dkf} one computes
    \begin{align}
    Df[\nu_1] = \sqrt{L}\sigma_1^{1-1/L}e_1(d_L)e_1(d_0)^T,
    \end{align}
    \begin{align}
    D^2f[\nu_1^{\odot 2}] = (L-1)\sigma_1^{1-2/L}e_1(d_L)e_1(d_0)^T,
    \end{align}
    and
    \begin{align}
    D^3f[\nu_1^{\odot 3}] =\frac{(L-2)(L-1)}{\sqrt{L}}
\sigma_1^{1-3/L}e_1(d_L)e_1(d_0)^T.
    \end{align}
    Using \eqref{eq:D4ell}, these identities imply that
    \begin{align}
        D^4\ell[\nu_1^{\odot 4}] &= 4\langle D^3f[\nu_1^{\odot 3}],Df[\nu_1]\rangle + 3\langle D^2f[\nu_1^{\odot 2}],D^2f[\nu_1^{\odot 2}]\rangle\\&=(L-1)(7L-11)\sigma_1^{2-4/L}.
    \end{align}
    To compute the term $D^3\ell[\nu_1^{\odot 2},A\nabla^3\ell[\nu_1^{\odot 2}]]$, using \eqref{eq:Dkf} and fixing a tangent vector $h = (h_1,\dots,h_L)\in\prod_{l=1}^L\RB^{d_l\times (d_{l-1})}$, one computes
    \begin{align}
    D^2f[\nu_1,h] = \frac{\sigma_1^{1-2/L}}{\sqrt{L}}\bigg(\sum_{l_1=1}^L\sum_{l_2\neq l_1}e_1(d_{l_2})^Th_{l_2}e_1(d_{l_2-1})\bigg) \,e_1(d_L)e_1(d_0)^T + B,
    \end{align}
    and
    \begin{align}
    Df[h] = \sigma_1^{1-1/L}\bigg(\sum_{l=1}^Le_1(d_l)^Th_le_1(d_{l-1})\bigg)e_1(d_L)e_1(d_0)^T + B'
    \end{align}
    where $B,B'\in \big(e_1(d_L)e_1(d_0)^T\big)^{\perp}$. Then by \eqref{eq:D3ell} one has
    \begin{align}
    D^3\ell[\nu_1^{\odot 2},h] &= 2\langle D^2f[\nu_1,h],Df[\nu_1]\rangle + \langle D^2f[\nu_1^{\odot 2}],Df[h]\rangle\\&=3(L-1)\sigma_1^{2-3/L}\sum_{l=1}^Le_1(d_l)^Th_le_1(d_{l-1})\\&=3\sqrt{L}(L-1)\sigma_1^{2-3/L}\langle\nu_1,h\rangle.
    \end{align}
    It follows that $\nabla^3\ell[\nu_1^{\odot 2}] = 3\sqrt{L}(L-1)\sigma_1^{2-3/L}\nu_1$, hence that
    \begin{align}
        A_{\eta}\nabla^3\ell[\nu_1^{\odot2}] = \frac{3\sqrt{L}(L-1)\sigma_1^{2-3/L}}{\lambda_1(1-\eta\lambda_1)}\nu_1,
    \end{align}
    implying further that
    \begin{align}
        D^3\ell[\nu_1^{\odot 2},A_{\eta}\nabla^3\ell[\nu_1^{\odot 2}]] &=\frac{3\sqrt{L}(L-1)\sigma_1^{2-3/L}}{\lambda_1(1-\eta\lambda_1)}D^3\ell[\nu_1^{\odot 3}]\\&=\frac{9L(L-1)^2\sigma_1^{4-6/L}}{\lambda_1(1-\eta\lambda_1)}\\&=\frac{9(L-1)^2\sigma_1^{2-4/L}}{1-\eta\lambda_1},
    \end{align}
    with the final line following from $\lambda_1|_{F} = L\sigma_1^{2-2/L}$. Putting everything together, one arrives at the claimed formula for $\alpha_{\eta}$.

    We now move on to proving that $D\alpha_{\eta}(w) = 0$. Since $\alpha_{\eta}|_{F}$ is constant, it suffices to show that $D\alpha_{\eta}(w)$ kills the normal directions to $F$ as worked out in Proposition \ref{prop:tangents}. Recall from Proposition \ref{prop:morsebott} that $\nu_wF$ splits into subspaces $V_1(w)$ and $V_2(w)$ consisting of modifications to the top-left entries of the factors and of rotations of their corresponding singular vectors respectively.
    
    It is relatively easy to show that $D\alpha_{\eta}(w)|_{V_2(w)}\equiv 0$. Recall the isometry $\Theta$ from \eqref{eq:Theta}, under which $f$ is equivariant and for which $D\Theta(w)$ has $-1$ eigenspace precisely equal to $V_2(w)$. Since  $f$ is equivariant under $\Theta$, $\alpha_{\eta}$ is invariant under $\Theta$; and since $\Theta(w) = w$, $D\alpha_{\eta}(w)$ is also invariant under $\Theta$ so that for any $\xi\in V_2$ one has
    \begin{align}
    D\alpha_{\eta}(w)[\xi] = D\alpha_{\eta}(w)[D\Theta(w)\xi] = -D\alpha_{\eta}(w)[\xi],
    \end{align}
    implying that $D\alpha_{\eta}(w)[\xi] = 0$ and hence that $D\alpha_{\eta}(w)|_{V_2(w)}\equiv 0$.

    We now turn to demonstrating $D\alpha_{\eta}(w)|_{V_1(w)}\equiv 0$. Fix a tuple $a = (a_l)_{l=1}^{L-1}\in\RB^{L-1}$ and for notational convenience denote $r_l:=a_l-a_{l-1}$ for all $l=1,\dots,L$, where we set $a_0,a_L:=0$. As in Proposition \ref{prop:morsebott}, we consider the curve
    \begin{align}
    \gamma:t\mapsto \Bigg(\begin{pmatrix}\exp(tr_l)\sigma_1^{1/L} & 0 \\ 0 & \bW_l\end{pmatrix}\Bigg)_{l=1}^L,
    \end{align}
    with $\gamma(0) = w$, and evaluate $\alpha_{\eta}(\gamma(t))$. To this end, we denote
    \begin{align}
    S_j(t):=\sum_{l=1}^L\exp(-2jtr_l),\qquad j=1,2,3.
    \end{align}
    One has
    \begin{align}
    \nu_1(\gamma(t)) &= \frac{1}{\|Df(\gamma(t))^T[e_1(d_L)e_1(d_0)^T]\|}Df(\gamma(t))^T[e_1(d_L)e_1(d_0)^T]\\&=S_1(t)^{-1/2}\big(\exp(-tr_l)e_1(d_l)e_1(d_{l-1})^T\big)_{l=1}^L.
    \end{align}
    From now on, we suppress evaluation at $\gamma(t)$ for notational convenience. Using the fact that $\prod_{l=1}^L\exp(tr_l) = 1$, note that for any tuple $l_1<\dots<l_k$ one has
    \begin{align}\label{eq:prod}
        \prod_{m\neq l_1,\dots,l_k}\exp(tr_m) = \prod_{j=1}^k\exp(-tr_{l_j}).
    \end{align}
    We now turn to computing the quantities appearing in \eqref{eq:alpha}. Invoking \eqref{eq:Dkf} gives
    \begin{align}
        Df[\nu_1] &= \bigg(\sum_{l=1}^LS_1^{-1/2}\exp(-tr_l)\sigma_1^{1-1/L}\prod_{m\neq l}\exp(tr_m)\bigg)e_1(d_L)e_1(d_0)^T\\&=S_1^{-1/2}\sigma_1^{1-1/L}\bigg(\sum_{l=1}^L\exp(-2tr_l)\bigg)e_1(d_L)e_1(d_0)^T\\&=\sigma_1^{1-1/L}S_1^{1/2}e_1(d_L)e_1(d_0)^T
    \end{align}
    by invoking \eqref{eq:prod} for the second line. Similarly, invoking \eqref{eq:Dkf} and \eqref{eq:prod} gives
    \begin{align}
        D^2f[\nu_1^{\odot 2}] &= \sigma_1^{1-2/L}S_1^{-1}\bigg(\sum_{l_1=1}^L\sum_{l_2\neq l_1}\exp(-tr_{l_1})\exp(-tr_{l_2})\prod_{m\neq l_1,l_2}\exp(tr_m)\bigg)e_1(d_L)e_1(d_0)^T\\&=\sigma_1^{1-2/L}S_1^{-1}\bigg(\sum_{l_1=1}^L\exp(-2tr_{l_1})\sum_{l_2\neq l_1}\exp(-2tr_{l_2})\bigg)e_1(d_L)e_1(d_0)^T\\&=\sigma_1^{1-2/L}\frac{S_1^2-S_2}{S_1}e_1(d_L)e_1(d_0)^T
    \end{align}
    by using the identity $\sum_{l_1\neq l_2}p_{l_1}p_{l_2} = \big(\sum_lp_l\big)^2 - \sum_lp_l^2$ for the final line, and
    \begin{align}
        D^3f[\nu_1^{\odot 3}] &= S_1^{-3/2}\sigma_1^{1-3/L}\bigg(\sum_{l_1=1}^L\sum_{l_2\neq l_1}\sum_{l_3\neq l_2,l_1}\exp(-t(r_{l_1} + r_{l_2} + r_{l_3}))\prod_{m\neq l_1,l_2,l_3}\exp(tr_m)\bigg)e_1(d_L)e_1(d_0)^T\\&=S_1^{-3/2}\sigma_1^{1-3/L}\bigg(\sum_{l_1\neq l_2\neq l_3}\exp(-2t(r_{l_1}+r_{l_2}+r_{l_3}))\bigg)e_1(d_L)e_1(d_0)^T\\&=\sigma_1^{1-3/L}\frac{S_1^3-3S_2S_1 + 2S_3}{S_1^{3/2}}e_1(d_L)e_1(d_0)^T
    \end{align}
    by using the identity $\sum_{l_1\neq l_2\neq l_3}p_{l_1}p_{l_2}p_{l_3} = \big(\sum_lp_l\big)^3-3\big(\sum_lp_l\big)\big(\sum_lp_l^2\big) + 2\big(\sum_lp_l^3\big)$ for the final line. It follows from \eqref{eq:D4ell} that
    \begin{align}
        D^4\ell[\nu_1^{\odot 4}] = \sigma_1^{2-4/L}\bigg(4\frac{S_1^3-3S_1S_2+2S_3}{S_1} + 3\bigg(\frac{S_1^2-S_2}{S_1}\bigg)^2\bigg).
    \end{align}
    Finally, we come to computing the term $D^3\ell[\nu_1^{\odot 2},A_{\eta}\nabla^3\ell[\nu_1^{\odot 2}]]$. For arbitrary $h = (h_1,\dots,h_L)\in\prod_{l=1}^L\RB^{d_l\times d_{l-1}}$, using \eqref{eq:Dkf} and \eqref{eq:prod} one has
    \begin{align}
        Df[h] &= \sigma_1^{1-1/L}\bigg(\sum_{l=1}^L\prod_{m\neq l}\exp(tr_m)\,e_1(d_l)^Th_le_1(d_{l-1})\bigg)e_1(d_L)e_1(d_0)^T + C\\&=\sigma_1^{1-1/L}\bigg(\sum_{l=1}^L\exp(-tr_l)\,e_1(d_l)^Th_le_1(d_{l-1})\bigg)e_1(d_L)e_1(d_0)^T + C
    \end{align}
    for some $C\in \big(e_1(d_L)e_1(d_0)\big)^{\perp}$, and
    \begin{align}
        D^2f[\nu_1,h] &= \sigma_1^{1-2/L}S_1^{-1/2}\bigg(\sum_{l_1=1}^L\sum_{l_2\neq l_1}\exp(-tr_{l_1})\,e_1(d_{l_2})^Th_{l_2}e_1(d_{l_2-1})\prod_{m\neq l_1,l_2}\exp(tr_m)\bigg)e_1(d_L)e_1(d_0)^T + C'\\&=\sigma_1^{1-2/L}S_1^{-1/2}\bigg(\sum_{l_1=1}^L\sum_{l_2\neq l_1}\exp(-t(2r_{l_1} + r_{l_2}))e_1(d_{l_2})^Th_{l_2}e_1(d_{l_2-1})\bigg)e_1(d_L)e_1(d_0)^T + C'
    \end{align}
    for some $C'\in\big(e_1(d_L)e_1(d_0)^T\big)^{\perp}$. By \eqref{eq:D3ell}, it follows that
    \begin{align}
        D^3\ell[\nu_1^{\odot 2},h] &= 2\langle D^2f[\nu_1,h],Df[\nu_1]\rangle + \langle D^2f[\nu_1^{\odot 2}],Df[h]\rangle\\&=\sigma_1^{2-3/L}\sum_{l=1}^L\bigg(\exp(-tr_l)\bigg(\frac{S_1^2-S_2}{S_1} + 2(S_1-\exp(-2tr_l))\bigg)e_1(d_l)^Th_le_1(d_{l-1})\bigg),
    \end{align}
    so that
    \begin{align}
        \nabla^3\ell[\nu_1^{\odot 2}] = \bigg(\sigma_1^{2-3/L}\exp(-tr_l)\bigg(\frac{S_1^2-S_2}{S_1} + 2(S_1-\exp(-2tr_l))\bigg)e_1(d_l)e_1(d_{l-1})^T\bigg)_{l=1}^L.
    \end{align}
    Since each of the components of this tuple is nonzero only in the top-left entry, its projection into the span of $\nu_{2:q}$ is zero while its projection into the span of $\nu_1$ is given by
    \begin{align}
        \nu_1^T\nabla^3\ell[\nu_1^{\odot 2}] &= D^3\ell[\nu_1^{\odot 3}] = 3\sigma_1^{2-3/L}\frac{S_1^2-S_2}{S_1^{1/2}}
    \end{align}
    by \eqref{eq:D3ell}, hence
    \begin{align}
        D^3\ell[\nu_1^{\odot 2},A\nabla^3\ell[\nu_1^{\odot 2}]] &= \lambda_1^{-1}(1-\eta\lambda_1)^{-1}\big(D^3\ell[\nu_1^{\odot 3}]\big)^2\\&=\frac{9\sigma_1^{2-4/L}(S_1^2-S_2)^2}{(1-\eta\lambda_1)S_1^2}
    \end{align}
    after substituting $\lambda_1 = \sigma_1^{2-2/L}S_1$. Thus, along the curve $\gamma(t)$ one has
    \begin{align}
        \alpha_{\eta} = -\sigma_1^{2-4/L}\bigg(\frac{9(S_1^2-S_2)^2}{(1-\eta\lambda_1)S_1^2} + \frac{4(S_1^3-3S_2S_1+2S_3)}{3S_1} + \frac{(S_1^2-S_2)^2}{S_1^2}\bigg).
    \end{align}
    Finally, note that for any $j=1,2,3$ one has $\dot{S}_j(0) = -2j\sum_{l=1}^Lr_l = -2j\sum_{l=1}^L(a_l-a_{l-1}) = 0$ since $a_0 = a_L = 0$, and similarly $\dot{\lambda_1}(0) = \sigma_1^{2-2/L}\dot{S}_1(0) = 0$. It follows that $\dot{\alpha}_{\eta}(0) = 0$, thus completing the proof.
\end{proof}

\end{document}
